\begin{introduction}

La Langue des Signes Malgache (LSM) est la langue utilisée par les personnes sourdes et leurs proches à Madagascar. Elle possède sa propre organisation linguistique et mobilise plusieurs parties du corps : les mains, le visage, les épaules et le regard. Dans le pays, les personnes sourdes accèdent difficilement aux services essentiels. Elles ont peu d'interlocuteurs capables de signer dans les hôpitaux, les écoles ou les administrations. Des outils de traduction adaptés au contexte local pourraient réduire considérablement cet écart.

Ce travail porte sur la LSM. Ce choix s'explique par la volonté de traiter une langue utilisée localement, encore peu représentée dans les travaux de reconnaissance automatique. Un dictionnaire de la LSM a été publié par la Fédération des Sourds de Madagascar, avec des entrées en malgache et leurs équivalents signés. Ce dictionnaire fournit une base lexicale utilisable pour construire un corpus de signes. La méthode proposée reste transposable à d'autres langues des signes dès qu'un vocabulaire équivalent est défini.

Les recherches récentes en reconnaissance de la langue des signes utilisent surtout des modèles profonds : réseaux convolutifs, Transformers, réseaux de convolution sur graphes. Ces modèles obtiennent de bons résultats sur de grands corpus. En revanche, ils demandent beaucoup de données annotées et un temps d'entraînement long. Une autre voie consiste à travailler sur les points du squelette extraits de la vidéo, plutôt que sur les images brutes. Cette représentation décrit directement la position du corps et des mains. Elle permet de comparer plusieurs modèles sur les mêmes entrées. C'est cette voie que suit le présent travail.

La question de recherche est la suivante : dans quelle mesure un modèle récurrent, entraîné sur des points du squelette extraits par MediaPipe Holistic, permet-il de traduire de façon fiable des signes isolés de la LSM ?

Le périmètre est volontairement restreint : dix signes, quatre signeurs dont un réservé au test, signes isolés, vocabulaire fermé et prototype local. Ces limites sont des choix méthodologiques. Elles garantissent la reproductibilité du pipeline et la traçabilité des décisions. La contribution visée consiste à documenter une démarche expérimentale et à fournir un prototype reproductible.

Le mémoire comporte cinq chapitres. Le premier présente l'état de l'art : spécificités de la LSM, représentations des signes, architectures de reconnaissance et travaux consacrés aux langues des signes. Le deuxième décrit la méthodologie : protocole d'enregistrement, extraction des points du squelette, préparation des données et protocole d'évaluation. Le troisième expose les expériences et les résultats de la comparaison entre LSTM et GRU, ainsi que le test de robustesse. Le quatrième présente la conception et la réalisation de l'application web, son interface et ses tests d'intégration. Le cinquième discute des limites du modèle et de l'application, puis ouvre sur les perspectives.

\end{introduction}
